% Compile with xelatex (twice).
% Tables come from 10_01_descriptive_tables.py via tables_naver.tex, which defines
% \tblVariables \tblDescriptives \tblCorrelations \tblVIF as reusable macros —
% the document is bilingual and prints each table once per language.
\documentclass[10pt,a4paper]{article}

\usepackage{fontspec}
\usepackage{xeCJK}
\setCJKmainfont{Malgun Gothic}
\setCJKsansfont{Malgun Gothic}
% Korean uses spaces between words, unlike Chinese/Japanese: keep them visible
\xeCJKsetup{CJKspace=true}

\usepackage[top=1.1cm, bottom=1.0cm, left=1.4cm, right=1.4cm]{geometry}
\usepackage{booktabs}
\usepackage{multirow}
\usepackage{tikz}
\usetikzlibrary{arrows.meta,positioning,fit,backgrounds}
\usepackage{amsmath}
\usepackage{amssymb}
\usepackage{enumitem}
\usepackage{microtype}
\usepackage{caption}
\usepackage{float}   % [H] placement used by the generated tables
\usepackage{url}
% Match the reference paper's caption style: "표 1. 변수요약 / Variable summary"
\captionsetup[table]{name=표, labelsep=period, font=small, labelfont=bf,
                     skip=3pt, justification=centering}

\linespread{0.96}
\setlength{\parindent}{0pt}
\setlength{\parskip}{2pt}
\setlength{\emergencystretch}{1em}
\renewcommand{\arraystretch}{1.0}
\pagestyle{plain}

\newcommand{\headrule}{\vspace{1pt}\noindent\rule{\linewidth}{0.4pt}\vspace{1pt}}
\newcommand{\sect}[1]{\vspace{3pt}\noindent\textbf{#1}\par\vspace{1pt}}

% Generated by 10_01_descriptive_tables.py -- do not edit by hand.
% Requires \usepackage{float} in the including document ([H] placement).
% Defines: \tblPlaces \tblVariables \tblDescriptives \tblCorrelations \tblVIF

\newcommand{\tblPlaces}{%
\begin{table}[H]\centering%
\caption{관광지별 리뷰 수 / Reviews per attraction}%
\footnotesize\begin{tabular}{llrrr}\hline%
관광지 & English & Reviews & w/ text & w/ rating \\ \hline%
서울스카이 & Seoul Sky & 1,092 & 1,025 & 792 \\%
남산골한옥마을 & Namsangol Village & 1,075 & 1,033 & 826 \\%
에버랜드 & Everland & 995 & 935 & 851 \\%
N서울타워 & N Seoul Tower & 939 & 817 & 198 \\%
동대문디자인플라자 & DDP & 668 & 453 & 232 \\%
덕수궁 & Deoksugung & 646 & 539 & 72 \\%
운현궁 & Unhyeongung & 616 & 562 & 384 \\%
서대문형무소역사관 & Seodaemun Prison History Hall & 610 & 331 & 27 \\%
창덕궁 & Changdeokgung & 610 & 399 & 38 \\%
올림픽공원 & Olympic Park & 609 & 499 & 104 \\%
해운대해수욕장 & Haeundae Beach & 608 & 457 & 53 \\%
서대문자연사박물관 & Seodaemun Museum & 608 & 320 & 20 \\%
광안리해수욕장 & Gwangalli Beach & 607 & 505 & 77 \\%
국립중앙박물관 & Natl. Museum & 606 & 543 & 221 \\%
여의도한강공원 & Yeouido Park & 605 & 443 & 87 \\%
코엑스 & COEX & 590 & 359 & 198 \\%
창경궁 & Changgyeonggung & 590 & 418 & 26 \\%
경복궁 & Gyeongbokgung & 580 & 477 & 91 \\%
낙산공원 & Naksan Park & 435 & 249 & 108 \\%
북촌한옥마을 & Bukchon Village & 296 & 165 & 118 \\%
북한산 & Bukhansan & 253 & 155 & 26 \\%
롯데월드 & Lotte World & 180 & 85 & 10 \\%
경희궁 & Gyeonghuigung & 170 & 98 & 12 \\%
은평한옥마을 & Eunpyeong Village & 156 & 125 & 36 \\%
인사동길 & Insadong-gil & 109 & 68 & 6 \\%
광장시장 & Gwangjang Mkt & 59 & 32 & 14 \\%
명동성당 & Myeongdong Cathedral & 39 & 33 & 18 \\%
홍대거리 & Hongdae St. & 30 & 22 & 6 \\%
서촌 & Seochon & 20 & 18 & 1 \\%
\hline\end{tabular}%
\end{table}%
}

\newcommand{\tblVariables}{%
\begin{table}[H]\centering%
\caption{변수요약 / Variable summary}%
\footnotesize\begin{tabular}{llp{7.6cm}}\hline%
변수 종류 & 변수 이름 & 변수 설명 \\ \hline%
종속변수 (DV) & Revisit (재방문) & 해당 관광지를 2회차 이상 방문한 작성자의 리뷰 여부 (0/1) \\%
 & Keyword Count (키워드 수) & 리뷰당 작성자가 선택한 키워드 수 (0$\sim$5), 정보 다양성 지표 \\%
 & Rating (별점) & 작성자가 부여한 별점 (1$\sim$5) \\%
 & Reactions (반응 수) & 해당 리뷰에 반응을 남긴 이용자 수, 유용성 지표 \\%
독립변수 (IV) & Companion (동반자) & 방문 동반자 유형 더미: 혼자·연인·친구·아이·가족·동료 \\%
 & Purpose (방문 목적) & 방문 목적 더미: 여행·나들이·일상·데이트·친목 \\%
 & Waited (대기) & 방문 시 대기 경험 여부 (0/1) \\%
 & Reserved (예약) & 사전 예약 여부 (0/1) \\%
 & Sentiment (감성 점수) & 리뷰 텍스트의 감성 점수 (-1 부정 $\sim$ 1 긍정) \\%
 & Tag themes (키워드 주제) & 키워드를 10개 주제로 분류한 더미: 시각·볼거리·혼잡·한적·시설·접근성·가족·가격·음식·물놀이 \\%
통제변수 (Control) & Text Length (텍스트 길이) & 리뷰 텍스트 길이의 로그값 \\%
 & Image Count (사진 수) & 리뷰에 첨부된 사진 수 \\%
 & User Reviews (작성자 리뷰 수) & 작성자가 누적 작성한 리뷰 수의 로그값 \\%
 & User Followers (팔로워 수) & 작성자의 팔로워 수의 로그값 \\%
 & Review Age (리뷰 경과일) & 리뷰 작성 후 경과일의 로그값, 반응 수 모형에서 노출 기간을 통제 \\%
 & Verification (인증) & 영수증 인증 리뷰 여부 (0/1) \\%
 & Season (계절) & 방문 계절 더미: 여름·가을·겨울 (봄 기준) \\%
 & Weekend (주말) & 주말 방문 여부 (0/1) \\%
 & Place FE (관광지 고정효과) & 관광지별 고정효과, 장소 고유 특성 통제 \\%
\hline\end{tabular}\end{table}%
}

\newcommand{\tblDescriptives}{%
\begin{table}[H]\centering%
\caption{기술통계표 / Descriptive statistics}%
\footnotesize\begin{tabular}{lrrrrrr}\hline%
Statistic & N & Mean & St.\ Dev. & Min & Median & Max \\ \hline%
Revisit (재방문) & 14,401 & 0.218 & 0.413 & 0 & 0 & 1 \\%
Keyword Count (키워드 수) & 14,401 & 2.97 & 1.883 & 0 & 3 & 5 \\%
Rating (별점) & 4,652 & 4.682 & 0.732 & 0.5 & 5 & 5 \\%
Reactions (반응 수) & 14,401 & 1.437 & 7.415 & 0 & 0 & 122 \\%
Sentiment (감성 점수) & 11,165 & 0.653 & 0.457 & -0.973 & 0.87 & 0.988 \\%
Text Length (텍스트 길이) & 11,165 & 112.089 & 106.805 & 14 & 71 & 753 \\%
Image Count (사진 수) & 14,401 & 0.235 & 1.412 & 0 & 0 & 20 \\%
User Reviews (작성자 리뷰 수) & 14,401 & 823.26 & 1,009.4 & 1 & 482 & 19,082.0 \\%
Review Age (리뷰 경과일) & 14,401 & 588.324 & 558.142 & 0 & 428 & 3,957.0 \\%
\hline\end{tabular}\end{table}%
}

\newcommand{\tblCorrelations}{%
\begin{table}[H]\centering%
\caption{상관분석표 / Correlation matrix}%
\footnotesize\begin{tabular}{lrrrrrrrrr}\hline%
 & 1 & 2 & 3 & 4 & 5 & 6 & 7 & 8 & 9 \\ \hline%
1. Revisit (재방문) & 1.000 &  &  &  &  &  &  &  &  \\%
2. Keyword Count (키워드 수) & 0.140 & 1.000 &  &  &  &  &  &  &  \\%
3. Rating (별점) & 0.069 & 0.264 & 1.000 &  &  &  &  &  &  \\%
4. Reactions (반응 수) & 0.194 & 0.097 & 0.051 & 1.000 &  &  &  &  &  \\%
5. Sentiment (감성 점수) & -0.067 & 0.116 & 0.514 & -0.039 & 1.000 &  &  &  &  \\%
6. Text Length (텍스트 길이) & 0.198 & 0.204 & 0.000 & 0.221 & -0.095 & 1.000 &  &  &  \\%
7. Image Count (사진 수) & 0.036 & 0.075 & 0.044 & 0.057 & 0.016 & 0.101 & 1.000 &  &  \\%
8. User Reviews (작성자 리뷰 수) & 0.255 & 0.031 & 0.036 & 0.212 & -0.082 & 0.159 & 0.046 & 1.000 &  \\%
9. Review Age (리뷰 경과일) & -0.140 & -0.395 & -0.094 & -0.083 & -0.027 & -0.137 & -0.057 & 0.034 & 1.000 \\%
\hline\end{tabular}\end{table}%
}

\newcommand{\tblVIF}{%
\begin{table}[H]\centering%
\caption{분산팽창계수 / Variance inflation factors}%
\footnotesize\begin{tabular}{lrlrlr}\hline%
변수 / Variable & VIF & 변수 / Variable & VIF & 변수 / Variable & VIF \\ \hline%
has\_visit\_info & 2.82 & comp\_family & 1.21 & sentiment & 1.06 \\%
waited & 1.10 & comp\_colleague & 1.25 & n\_tags & 1.54 \\%
reserved & 1.84 & purp\_travel & 1.38 & revisit & 1.16 \\%
comp\_alone & 1.69 & purp\_outing & 1.65 & ln\_length & 1.10 \\%
comp\_partner & 1.75 & purp\_daily & 1.45 & ln\_user\_reviews & 1.54 \\%
comp\_friend & 1.53 & purp\_date & 1.67 &  &  \\%
comp\_child & 1.70 & purp\_social & 1.39 &  &  \\%
\hline\end{tabular}%
\par\vspace{2pt}\footnotesize\textit{All values below 5; multicollinearity is not a concern. 모든 값이 5 미만으로 다중공선성 문제는 없다.}%
\end{table}%
}

% Shared body for the variable-map table; used by \tblVariableMapEN and \tblVariableMapKO.
% Notation key: $Y_{ij}$ = outcome of that model; $\beta_k$ = coefficient in x_ij vector;
%   $\delta_j$ = place fixed effect; $\gamma$ = ln_review_age exposure offset (M4 only);
%   $\beta_k{\times}N$ = N binary dummies each with their own \beta_k; -- = not included.
% Scale abbreviations used below:
%   bin. = binary (0/1)   ord. = ordinal   cont. = continuous   count = count (integer)   categ. = categorical
\newcommand{\tblVariableMapEN}{%
\begin{table}[H]\centering%
\caption{Variable-to-model mapping}%
\footnotesize\setlength{\tabcolsep}{3.5pt}%
\begin{tabular}{lp{2.7cm}p{2.8cm}lp{1.4cm}p{1.4cm}p{1.4cm}p{1.4cm}}\hline%
Type & Concept (Table~2) & Dataset column(s) & Scale & M1 & M2 & M3 & M4 \\\hline%
\multicolumn{8}{l}{\quad\textit{Dependent variables}}\\\hline%
DV & Revisit             & \texttt{revisit}                          & bin.          & $Y_{ij}$              & $\beta_k$             & $\beta_k$             & $\beta_k$             \\%
DV & Keyword Count       & \texttt{n\_tags}                          & ord.\ 0--5    & $\beta_k$             & $Y_{ij}$              & $\beta_k$             & $\beta_k$             \\%
DV & Rating              & \texttt{rating}                           & ord.\ 1--5    & --                    & --                    & $Y_{ij}$              & --                    \\%
DV & Reactions           & \texttt{reactions}                        & count         & --                    & --                    & --                    & $Y_{ij}$              \\\hline%
\multicolumn{8}{l}{\quad\textit{Independent variables}}\\\hline%
IV & Companion           & \texttt{comp\_*} ($\times$6)              & bin.          & $\beta_k{\times}6$    & $\beta_k{\times}6$    & $\beta_k{\times}4^a$  & $\beta_k{\times}6$    \\%
IV & Purpose             & \texttt{purp\_*} ($\times$5)              & bin.          & $\beta_k{\times}5$    & $\beta_k{\times}5$    & $\beta_k{\times}2^a$  & $\beta_k{\times}5$    \\%
IV & Waited              & \texttt{waited}                           & bin.          & $\beta_k$             & $\beta_k$             & --                    & $\beta_k$             \\%
IV & Reserved            & \texttt{reserved}                         & bin.          & $\beta_k$             & $\beta_k$             & $\beta_k$             & $\beta_k$             \\%
IV & Sentiment           & \texttt{sentiment}                        & cont.         & $\beta_k$             & --                    & --$^b$                & $\beta_k$             \\%
IV & Tag themes          & \texttt{tag\_*} ($\times$10)              & bin.          & --                    & --                    & --                    & --$^c$                \\%
IV & Has text            & \texttt{has\_text}                        & bin.          & --                    & $\beta_k$             & $\beta_k$             & --                    \\%
IV & Visit info          & \texttt{has\_visit\_info}                 & bin.          & $\beta_k$             & $\beta_k$             & --                    & $\beta_k$             \\\hline%
\multicolumn{8}{l}{\quad\textit{Control variables}}\\\hline%
Ctrl & Text Length       & \texttt{ln\_length}                       & cont.         & $\beta_k$             & --                    & $\beta_k$             & $\beta_k$             \\%
Ctrl & User Reviews      & \texttt{ln\_user\_reviews}                & cont.         & $\beta_k$             & $\beta_k$             & $\beta_k$             & $\beta_k$             \\%
Ctrl & User Followers    & \texttt{ln\_user\_followers}              & cont.         & --                    & --                    & --                    & $\beta_k$             \\%
Ctrl & Review Age        & \texttt{ln\_review\_age}                  & cont.         & --                    & --                    & --                    & $\gamma$              \\%
Ctrl & Season            & \texttt{season\_*} ($\times$3)            & bin.          & --                    & --                    & --                    & $\beta_k{\times}3$    \\%
Ctrl & Weekend           & \texttt{visit\_weekend}                   & bin.          & --                    & --                    & --                    & $\beta_k$             \\%
Ctrl & Verification      & \texttt{verif\_receipt}                   & bin.          & --                    & --                    & --                    & $\beta_k$             \\%
Ctrl & Place FE          & \textit{(absorbed)}                       & categ.        & $\delta_j$            & $\delta_j$            & $\delta_j$            & $\delta_j$            \\\hline%
\end{tabular}%
\par\vspace{3pt}\footnotesize%
\textit{Scale:} bin.\ = binary (0/1); ord.\ = ordinal; cont.\ = continuous (log-transformed); count = non-negative integer.\quad
\textit{Notation link:}
$\beta_k$ = enters $\mathbf{x}_{ij}^\top\boldsymbol\beta$ with coefficient $\beta_k$;\quad
$Y_{ij}$ = outcome (left-hand side; not in $\mathbf{x}_{ij}$);\quad
$\delta_j$ = place fixed effect;\quad
$\gamma$ = \texttt{ln\_review\_age} as exposure offset in M4 (not a $\beta_k$);\quad
$\beta_k{\times}N$ = $N$ dummies each with their own $\beta_k$.\\%
$^a$~Partial dummies in M3: only 32\% of reviews have a rating.~%
$^b$~Excluded from M3: $r{=}0.51$ with rating; VIF risk.~%
$^c$~Tag themes only in extended models M1b, M1x.%
\end{table}%
}

\newcommand{\tblVariableMapKO}{%
\begin{table}[H]\centering%
\caption{변수 매핑표}%
\footnotesize\setlength{\tabcolsep}{3.5pt}%
\begin{tabular}{lp{2.7cm}p{2.8cm}lp{1.4cm}p{1.4cm}p{1.4cm}p{1.4cm}}\hline%
유형 & 변수 이름 (표~2) & 데이터셋 열 & 척도 & M1 & M2 & M3 & M4 \\\hline%
\multicolumn{8}{l}{\quad\textit{종속변수}}\\\hline%
종속 & 재방문              & \texttt{revisit}                          & 이진          & $Y_{ij}$              & $\beta_k$             & $\beta_k$             & $\beta_k$             \\%
종속 & 키워드 수           & \texttt{n\_tags}                          & 순서 0--5     & $\beta_k$             & $Y_{ij}$              & $\beta_k$             & $\beta_k$             \\%
종속 & 별점               & \texttt{rating}                           & 순서 1--5     & --                    & --                    & $Y_{ij}$              & --                    \\%
종속 & 반응 수            & \texttt{reactions}                        & 계수          & --                    & --                    & --                    & $Y_{ij}$              \\\hline%
\multicolumn{8}{l}{\quad\textit{독립변수}}\\\hline%
독립 & 동반자 유형         & \texttt{comp\_*} ($\times$6)              & 이진          & $\beta_k{\times}6$    & $\beta_k{\times}6$    & $\beta_k{\times}4^a$  & $\beta_k{\times}6$    \\%
독립 & 방문 목적           & \texttt{purp\_*} ($\times$5)              & 이진          & $\beta_k{\times}5$    & $\beta_k{\times}5$    & $\beta_k{\times}2^a$  & $\beta_k{\times}5$    \\%
독립 & 대기 경험           & \texttt{waited}                           & 이진          & $\beta_k$             & $\beta_k$             & --                    & $\beta_k$             \\%
독립 & 사전 예약           & \texttt{reserved}                         & 이진          & $\beta_k$             & $\beta_k$             & $\beta_k$             & $\beta_k$             \\%
독립 & 감성 점수           & \texttt{sentiment}                        & 연속          & $\beta_k$             & --                    & --$^b$                & $\beta_k$             \\%
독립 & 키워드 주제         & \texttt{tag\_*} ($\times$10)              & 이진          & --                    & --                    & --                    & --$^c$                \\%
독립 & 텍스트 입력 여부    & \texttt{has\_text}                        & 이진          & --                    & $\beta_k$             & $\beta_k$             & --                    \\%
독립 & 방문 정보 제공      & \texttt{has\_visit\_info}                 & 이진          & $\beta_k$             & $\beta_k$             & --                    & $\beta_k$             \\\hline%
\multicolumn{8}{l}{\quad\textit{통제변수}}\\\hline%
통제 & 텍스트 길이         & \texttt{ln\_length}                       & 연속          & $\beta_k$             & --                    & $\beta_k$             & $\beta_k$             \\%
통제 & 작성자 리뷰 수      & \texttt{ln\_user\_reviews}                & 연속          & $\beta_k$             & $\beta_k$             & $\beta_k$             & $\beta_k$             \\%
통제 & 팔로워 수           & \texttt{ln\_user\_followers}              & 연속          & --                    & --                    & --                    & $\beta_k$             \\%
통제 & 리뷰 경과일         & \texttt{ln\_review\_age}                  & 연속          & --                    & --                    & --                    & $\gamma$              \\%
통제 & 방문 계절           & \texttt{season\_*} ($\times$3)            & 이진          & --                    & --                    & --                    & $\beta_k{\times}3$    \\%
통제 & 주말 방문           & \texttt{visit\_weekend}                   & 이진          & --                    & --                    & --                    & $\beta_k$             \\%
통제 & 인증 (영수증)       & \texttt{verif\_receipt}                   & 이진          & --                    & --                    & --                    & $\beta_k$             \\%
통제 & 관광지 고정효과     & \textit{(흡수됨)}                         & 범주          & $\delta_j$            & $\delta_j$            & $\delta_j$            & $\delta_j$            \\\hline%
\end{tabular}%
\par\vspace{3pt}\footnotesize%
\textit{척도:} 이진 = 0/1 더미; 순서 = 순서형 정수; 연속 = 로그 변환 연속형; 계수 = 비음 정수.\quad
\textit{기호 연결:}
$\beta_k$ = $\mathbf{x}_{ij}^\top\boldsymbol\beta$ 벡터에 계수 $\beta_k$로 포함;\quad
$Y_{ij}$ = 종속변수(좌변; $\mathbf{x}_{ij}$에 미포함);\quad
$\delta_j$ = 관광지 고정효과;\quad
$\gamma$ = M4에서 노출 오프셋($\beta_k$ 아님);\quad
$\beta_k{\times}N$ = 더미 $N$개가 각각 고유한 계수로 포함.\\%
$^a$~M3 일부 더미: 별점 보유 리뷰 32\%에 불과.~%
$^b$~M3 제외: 별점과 $r{=}0.51$, 다중공선성 우려.~%
$^c$~키워드 주제는 확장 모형 M1b, M1x에서만.%
\end{table}%
}


\begin{document}
\footnotesize

% ═══════════════════════════════════════════════════════════════════
% ENGLISH
% ═══════════════════════════════════════════════════════════════════
\begin{center}
{\large\textbf{Is a Review Useful?}}\\[2pt]
{\large\textbf{Visit Context and Information Diversity in Korean Attraction Reviews}}\\[5pt]
{\normalsize Jorge Adan Soriano Hernandez}\\
{\small Department of Data Science, Hoseo University \quad$\cdot$\quad Research proposal, October 2026}
\end{center}

\headrule

\sect{1. Topic}
When deciding whether to visit or revisit a place, one of the strongest considerations travellers
weigh is checking its reviews, comments, reactions and photos, a direct expression of social media's
role in decision-making. Within this context, information diversity, that is, how distinct the words
across a set of reviews actually are, is understood to drive a place's performance independently of
review volume. Lee, Jeong and Kim (2026) established this at the business-week level on Naver Smart
Place data: the \textbf{diversity} of information in a review (how many distinct keywords it carries)
drives business performance separately from the sheer \textbf{quantity} of reviews. This study proposes
analysing points of interest across Korea through individual reviews and their diversity, taking the
same construct down to the level of the review itself and asking where that diversity comes from.

\sect{2. Problem}
A place's performance can be measured in several ways, yet reviews themselves rarely separate how
much was written from how informative it was. Lee et al.\ (2026) show that information diversity
predicts business performance independently of review quantity, but they measure performance through
platform-side traffic (page views, total clicks) aggregated by business and week, so the content of
the review itself, and the context in which it was written, stay unobserved. Their information
diversity is a weekly average; the companion, the purpose, the waiting experience and the advance
reservation behind each visit are never used, although Naver collects all four. What a review
actually discusses---the aspects of the visit it foregrounds, beyond its overall sentiment---is
likewise unobserved. The antecedents of an informative review are therefore unknown: we know
diversity matters, not who supplies it or what it is about.

\sect{3. Research questions}
\begin{enumerate}[leftmargin=14pt, itemsep=0pt, topsep=1pt]
\item Which visit contexts are associated with returning to the same attraction?
\item Which reviewers and visit contexts produce more diverse information, i.e.\ more keywords?
\item Which reviews do other users respond to, and does sentiment help or hurt?
\item Which aspects of the visit do different reviewer groups discuss, and does what a review
      talks about predict reader engagement beyond its overall sentiment?
\end{enumerate}

\sect{4. Data}
Naver Place reviews for Korean attractions, collected by scraping each place individually.
14{,}401 reviews from 8{,}276 reviewers across 29 attractions; 11{,}165 carry text and 4{,}652 carry a
star rating. The unit of analysis is the review $i$ at place $j$, not the business--week. Every
place in the catalogue (\texttt{places.json}) with a Naver id has been scraped; coverage ranges from
1{,}092 reviews (서울스카이) down to 20 (서촌).

\tblPlaces




\clearpage
\sect{5. Variables}
Four dependent variables, the visit-context regressors Naver records, and the controls that
separate a behavioural effect from an artefact of who wrote the review and when.


\tblVariables

\begin{figure}[H]\centering
\begin{tikzpicture}[
  every node/.style={font=\footnotesize},
  mb/.style={draw, rounded corners=4pt, fill=white, inner sep=10pt,
             text width=13.5cm, align=left},
  arr/.style={-{Stealth[length=6pt]}, thick},
]
\node[mb] (data) {
  \textbf{Data Acquisition / 데이터 수집}\\[4pt]
  \quad Naver Place scraper $\;\longrightarrow\;$
  29 attractions $\;\cdot\;$ 14{,}401 reviews $\;\cdot\;$ 8{,}276 reviewers
};
\node[mb, below=1.0cm of data] (varbox) {
  \textbf{Variables / 변수}\\[6pt]
  \quad\begin{tabular}{|p{5.4cm}|p{5.4cm}|}
  \hline
  \rule{0pt}{9pt}\textbf{Visit Context} & \textbf{Reviewer Profile}\\[2pt]
  $\cdot$ Revisit (binary) & $\cdot$ Followers\\
  $\cdot$ Visited alone    & $\cdot$ Review count\\
  $\cdot$ Waited           & $\cdot$ Photo count\\
  $\cdot$ Group type       & \\[3pt]
  \hline
  \end{tabular}
};
\node[mb, below=1.0cm of varbox] (modbox) {
  \textbf{Models / 모형}\\[6pt]
  \begin{center}\begin{tabular}{|c|c|c|c|}
  \hline
  \rule{0pt}{14pt}\quad\textbf{M1}\quad & \quad\textbf{M2}\quad
    & \quad\textbf{M3}\quad & \quad\textbf{M4}\quad\\[4pt]
  \hline
  \end{tabular}\end{center}
};
\draw[arr] (data.south) -- (varbox.north);
\draw[arr] (varbox.south) -- (modbox.north);
\end{tikzpicture}
\caption{Research framework / 연구 프레임워크}
\end{figure}

\sect{6. Descriptive statistics}

\tblDescriptives

\sect{7. Correlations}
No pair is strong enough to threaten the regressions. The largest, Rating with Sentiment
($0.514$), confirms the sentiment classifier tracks the stars reviewers actually gave.
Keyword Count with Review Age ($-0.395$) shows keyword use has grown over the window,
which is why review age is controlled rather than ignored.


\clearpage

\tblCorrelations

\sect{8. Models}
Review-level regressions with place fixed effects $\delta_j$ and standard errors clustered by
reviewer, since the same person often writes several reviews.
\vspace{-6pt}
\begin{align*}
\text{M1:}\quad & \operatorname{logit} P(\textit{Revisit}_{ij}) = \beta_0 + \mathbf{x}_{ij}^\top\boldsymbol\beta + \delta_j \\[-1pt]
\text{M2:}\quad & P(\textit{KeywordCount}_{ij} \le k) = \big[1 + e^{-(\tau_k - \mathbf{x}_{ij}^\top\boldsymbol\beta - \delta_j)}\big]^{-1}, \quad k = 0,\dots,4 \\[-1pt]
\text{M3:}\quad & P(\textit{Rating}_{ij} \le k) = \big[1 + e^{-(\tau_k - \mathbf{x}_{ij}^\top\boldsymbol\beta - \delta_j)}\big]^{-1}, \quad k = 1,\dots,4 \\[-1pt]
\text{M4:}\quad & \textit{Reactions}_{ij} \sim \text{NegBin}(\mu_{ij}), \quad
                  \ln \mu_{ij} = \beta_0 + \mathbf{x}_{ij}^\top\boldsymbol\beta + \gamma \ln(\textit{Age}_{ij}) + \delta_j
\end{align*}
\vspace{-9pt}

The link function follows the measurement scale of each outcome: binary for M1, ordinal for M2 and
M3, and an overdispersed count for M4 (variance $\approx 37\times$ the mean; Poisson ruled out).
Each model is re-estimated on one random review per reviewer and per place; only effects that
survive are reported as robust. Cross-validation is grouped by reviewer.

\begin{table}[H]\centering
\caption{Model notation / Notation summary}
\footnotesize
\begin{tabular}{llp{8.2cm}}\toprule
\textbf{Model} & \textbf{Symbol} & \textbf{Description} \\ \midrule
\multirow{4}{*}{M1}
  & $i,\,j$                            & Review $i$ at place $j$ \\
  & $\beta_0$                          & Intercept \\
  & $\mathbf{x}_{ij}^\top\boldsymbol\beta$ & Vector of regressors (visit context + controls) and their coefficients \\
  & $\delta_j$                         & Place fixed effect (one per attraction) \\ \midrule
\multirow{6}{*}{M2, M3}
  & $i,\,j$                            & Review $i$ at place $j$ \\
  & $\beta_0$                          & Intercept \\
  & $\mathbf{x}_{ij}^\top\boldsymbol\beta$ & Vector of regressors and their coefficients \\
  & $\delta_j$                         & Place fixed effect \\
  & $\tau_k$                           & Ordered logit threshold for category $k$ \\
  & $k$                                & Outcome category (M2: $0\!\dots\!4$;\; M3: $1\!\dots\!4$) \\ \midrule
\multirow{7}{*}{M4}
  & $i,\,j$                            & Review $i$ at place $j$ \\
  & $\beta_0$                          & Intercept \\
  & $\mathbf{x}_{ij}^\top\boldsymbol\beta$ & Vector of regressors and their coefficients \\
  & $\delta_j$                         & Place fixed effect \\
  & $\mu_{ij}$                         & Expected reaction count \\
  & $\gamma$                           & Log-age elasticity (offset for review age) \\
  & $\alpha$                           & Negative binomial dispersion ($\hat\alpha = 2.44$; variance $\approx 37\times$ mean) \\ \bottomrule
\end{tabular}
\end{table}

\tblVariableMapEN

\tblVIF

\sect{9. Results}
\vspace{-3pt}
\begin{table}[H]\centering
\caption{Results summary / Leading effects (odds ratios and incidence rate ratios)}
\footnotesize
\begin{tabular}{lllr}\toprule
\textbf{Model} & \textbf{Fit} & \textbf{Variable} & \textbf{OR / IRR} \\ \midrule
\multirow{5}{*}{M1 Revisit} & \multirow{5}{*}{\shortstack[l]{AUC 0.773\\(ext.\ 0.785)}}
  & Reviewer's total reviews & 1.52*** \\
  & & Visited alone            & 1.68**  \\
  & & Waited                   & 0.51*** \\
  & & Sentiment                & 0.70*** \\
  & & Travel group             & 0.56*** \\ \midrule
\multirow{4}{*}{M2 Keyword Count} & \multirow{4}{*}{\shortstack[l]{RMSE 1.37\\vs 1.88 base}}
  & Wrote text               & 4.02*** \\
  & & Visit info provided      & 2.53*** \\
  & & Revisitor                & 1.60*** \\
  & & Travel group             & 1.73*** \\ \midrule
\multirow{2}{*}{M3 Rating} & \multirow{2}{*}{exploratory}
  & Keyword count            & 1.56*** \\
  & & Revisitor                & 1.37*   \\ \midrule
\multirow{3}{*}{M4 Reactions} & \multirow{3}{*}{AUC 0.815}
  & Followers                & 1.97*** \\
  & & Text length              & 1.64*** \\
  & & Sentiment                & 0.77**  \\ \bottomrule
\end{tabular}
\end{table}
{\scriptsize * $p<0.05$, ** $p<0.01$, *** $p<0.001$. Ratios $<1$ indicate a negative association. OR/IRR $= e^{\hat\beta}$.}

An aspect-level reading of 11{,}164 reviews, labelled on six content dimensions (historical/cultural,
accessibility, crowding, food, visual, service/experience) with an LLM-based classifier, sharpens
the picture. The six aspects were specified a priori from tourism-literature grounds and corroborated
against unsupervised topic extraction (LDA at optimal $K=6$ and BERTopic $K=15$) prior to LLM
classification. Visits with children raise mentions in three aspects (historical, accessibility and
service: ORs 1.22--1.55). Service friction splits two ways: queueing more than doubles crowding
mentions (OR $2.13$), while advance reservation cuts them (OR $0.69$) and more than doubles service
mentions (OR $2.39$). Date-purpose visits raise food mentions by 51\%. On engagement, no aspect
mention reaches $p<.05$ after reviewer controls: on Naver Place, reader reactions are driven by
reviewer identity (followers, revisitor status, text length), not by what the review discusses.

\sect{10. References}
\begin{enumerate}\small
  \item Bae, E., Chang, M., Park, E., \& Kim, D. (2017). The effect of Hallyu on tourism in Korea. \textit{Journal of Open Innovation: Technology, Market, and Complexity}, 3, 22.
  \item Chang, S. (2021). Food tourism in Korea. \textit{Journal of Vacation Marketing}.
  \item Ko, S., Kang, S., Kang, H., \& Lee, M.~J. (2018). An exploration of foreign tourists' perceptions of Korean food tour. \textit{Asia Pacific Journal of Tourism Research}, 23(8), 833--846.
  \item Karabila, I., Darraz, N., EL-Ansari, A., Alami, N., \& EL~Mallahi, M. (2024). BERT-enhanced sentiment analysis for personalized e-commerce recommendations. \textit{Multimedia Tools and Applications}, 83, 56463--56488.
  \item Lee, M., Jeong, D., \& Kim, A.~J. (2026). The impact of information diversity on online reviews and business performance. \textit{Journal of Consumer Studies}, 37(1). \url{https://doi.org/10.35736/JCS.37.1.1}
  \item Wang, S., Karmakar, S., Wang, F., \& Pei, Y. (2025). Content dissimilarity and online review helpfulness: Contextual insights. \textit{Journal of Business Research}, 187, 115068.
\end{enumerate}

\clearpage
\setcounter{table}{0}   % the Korean half repeats 표 1-4, it does not continue them

% ═══════════════════════════════════════════════════════════════════
% KOREAN
% ═══════════════════════════════════════════════════════════════════
\begin{center}
{\large\textbf{리뷰는 유용한가?:}}\\[2pt]
{\large\textbf{한국 관광지 리뷰의 방문 맥락과 정보 다양성 연구}}\\[5pt]
{\normalsize 소리아노 에르난데스 호르헤 아단}\\
{\small 호서대학교 데이터사이언스학과 \quad$\cdot$\quad 연구 계획서, 2026년 10월}
\end{center}



\headrule

\sect{1. 주제}
장소를 방문하거나 재방문할지 결정할 때, 여행자가 가장 중요하게 고려하는 것 중 하나는 리뷰·댓글·반응·
사진을 확인하는 일이며, 이는 소셜 미디어가 의사결정에 미치는 영향을 직접적으로 보여준다. 이러한
맥락에서 정보 다양성, 즉 리뷰에 사용된 단어가 서로 얼마나 다른가는 리뷰 양과 무관하게 장소의 성과를
이끄는 요인으로 이해된다. 이민영·정동영·김지영(2026)은 이를 업체·주 단위에서 네이버 스마트 플레이스
자료로 입증하였다. 즉 리뷰에 담긴 정보의 \textbf{다양성}(서로 다른 키워드의 수)은 정보의 \textbf{양}과는
별개로 업체 성과에 영향을 미친다. 본 연구는 한국 전역의 관심 지점을 개별 리뷰와 그 다양성을 통해
분석하며, 같은 개념을 리뷰 자체의 수준으로 내려 그 다양성이 어디에서 비롯되는지를 묻는다.

\sect{2. 문제 제기}
장소의 성과는 여러 방식으로 측정될 수 있지만, 리뷰 자체는 얼마나 많이 썼는가와 얼마나 유익하게
썼는가를 거의 구분하지 않는다. 이민영 외(2026)는 정보 다양성이 리뷰 양과 무관하게 업체 성과를
예측함을 보였으나, 성과를 업체·주 단위로 집계된 플랫폼 측 트래픽(페이지 뷰, 총 클릭 수)으로 측정하여
리뷰 자체의 내용과 그 작성 맥락은 관측하지 못한다. 해당 연구의 정보 다양성은 주간 평균값이며, 그
배후에 있는 동반자, 방문 목적, 대기 경험은 네이버가 모두 수집하고 있음에도 활용되지 않는다. 따라서
유익한 리뷰의 선행 요인은 아직 밝혀지지 않았다. 다양성이 중요하다는 점은 알지만, 누가 그것을
제공하는지는 모른다.

\sect{3. 연구 질문}
\begin{enumerate}[leftmargin=14pt, itemsep=0pt, topsep=1pt]
\item 어떤 방문 맥락이 동일 관광지로의 재방문과 연관되는가?
\item 어떤 작성자와 방문 맥락이 더 다양한 정보, 즉 더 많은 키워드를 제공하는가?
\item 어떤 리뷰가 다른 이용자의 반응을 얻으며, 감성은 이에 도움이 되는가 아니면 저해하는가?
\item 어떤 방문 측면을 작성자 그룹별로 언급하며, 리뷰가 다루는 내용이 전반 감성 이상의
      이용자 반응을 예측하는가?
\end{enumerate}

\sect{4. 데이터}
한국 관광지에 대한 네이버 플레이스 리뷰를 관광지별로 개별 수집하였다. 관광지 29곳, 작성자 8{,}276명,
리뷰 14{,}401건이며 이 중 11{,}165건은 텍스트를, 4{,}652건은 별점을 포함한다. 분석 단위는 업체·주가 아니라
관광지 $j$에 대한 개별 리뷰 $i$이다. 카탈로그(\texttt{places.json})에 네이버 ID가 있는 관광지는 모두
수집되었으며, 관광지별 리뷰 수는 최대 1{,}092건(서울스카이)에서 최소 20건(서촌)까지 분포한다.

\tblPlaces

\sect{5. 변수}
네 개의 종속변수와, 네이버가 기록하는 방문 맥락 변수, 그리고 누가 언제 작성했는지에서 비롯되는 효과를
행동적 효과와 분리하기 위한 통제변수로 구성된다.

\tblVariables

\begin{figure}[H]\centering
\begin{tikzpicture}[
  every node/.style={font=\footnotesize},
  mb/.style={draw, rounded corners=4pt, fill=white, inner sep=10pt,
             text width=13.5cm, align=left},
  arr/.style={-{Stealth[length=6pt]}, thick},
]
\node[mb] (data) {
  \textbf{데이터 수집 / Data Acquisition}\\[4pt]
  \quad 네이버 플레이스 스크레이퍼 $\;\longrightarrow\;$
  관광지 29곳 $\;\cdot\;$ 리뷰 14{,}401건 $\;\cdot\;$ 작성자 8{,}276명
};
\node[mb, below=1.0cm of data] (varbox) {
  \textbf{변수 / Variables}\\[6pt]
  \quad\begin{tabular}{|p{5.4cm}|p{5.4cm}|}
  \hline
  \rule{0pt}{9pt}\textbf{방문 맥락} & \textbf{작성자 프로필}\\[2pt]
  $\cdot$ 재방문 여부 & $\cdot$ 팔로워 수\\
  $\cdot$ 혼자 방문   & $\cdot$ 리뷰 수\\
  $\cdot$ 대기 경험   & $\cdot$ 사진 수\\
  $\cdot$ 동행 유형   & \\[3pt]
  \hline
  \end{tabular}
};
\node[mb, below=1.0cm of varbox] (modbox) {
  \textbf{모형 / Models}\\[6pt]
  \begin{center}\begin{tabular}{|c|c|c|c|}
  \hline
  \rule{0pt}{14pt}\quad\textbf{M1}\quad & \quad\textbf{M2}\quad
    & \quad\textbf{M3}\quad & \quad\textbf{M4}\quad\\[4pt]
  \hline
  \end{tabular}\end{center}
};
\draw[arr] (data.south) -- (varbox.north);
\draw[arr] (varbox.south) -- (modbox.north);
\end{tikzpicture}
\caption{연구 프레임워크 / Research framework}
\end{figure}

\sect{6. 기술통계}

\tblDescriptives

\sect{7. 상관분석}
회귀분석에 문제가 될 만큼 강한 상관은 없다. 가장 큰 값인 별점과 감성 점수의 상관($0.514$)은
감성 분류기가 실제 별점을 잘 따라가고 있음을 확인해 준다. 키워드 수와 리뷰 경과일의 상관($-0.395$)은
관측 기간 동안 키워드 사용이 증가했음을 의미하며, 이것이 리뷰 경과일을 반드시 통제해야 하는 이유이다.

\tblCorrelations

\sect{8. 연구모형}
관광지 고정효과 $\delta_j$를 포함한 리뷰 단위 회귀분석이며, 동일 작성자가 여러 리뷰를 작성하므로
표준오차는 작성자 단위로 군집화한다.
\vspace{-6pt}
\begin{align*}
\text{M1:}\quad & \operatorname{logit} P(\text{재방문}_{ij}) = \beta_0 + \mathbf{x}_{ij}^\top\boldsymbol\beta + \delta_j \\[-1pt]
\text{M2:}\quad & P(\text{키워드 수}_{ij} \le k) = \big[1 + e^{-(\tau_k - \mathbf{x}_{ij}^\top\boldsymbol\beta - \delta_j)}\big]^{-1}, \quad k = 0,\dots,4 \\[-1pt]
\text{M3:}\quad & P(\text{별점}_{ij} \le k) = \big[1 + e^{-(\tau_k - \mathbf{x}_{ij}^\top\boldsymbol\beta - \delta_j)}\big]^{-1}, \quad k = 1,\dots,4 \\[-1pt]
\text{M4:}\quad & \text{반응 수}_{ij} \sim \text{NegBin}(\mu_{ij}), \quad
                  \ln \mu_{ij} = \beta_0 + \mathbf{x}_{ij}^\top\boldsymbol\beta + \gamma \ln(\text{경과일}_{ij}) + \delta_j
\end{align*}
\vspace{-9pt}

연결함수는 각 종속변수의 측정 척도에 따라 결정된다. M1은 이항, M2와 M3은 순서형, M4는 과대산포된
계수형이다(분산 $\approx$ 평균의 37배; 포아송 기각). 강건성 검증을 통과한 효과만 보고하며,
교차검증은 작성자 단위로 분할한다.

\begin{table}[H]\centering
\caption{모형 기호 설명}
\footnotesize
\begin{tabular}{llp{8.2cm}}\toprule
\textbf{모형} & \textbf{기호} & \textbf{설명} \\ \midrule
\multirow{4}{*}{M1}
  & $i,\,j$                            & 관광지 $j$의 리뷰 $i$ \\
  & $\beta_0$                          & 절편 \\
  & $\mathbf{x}_{ij}^\top\boldsymbol\beta$ & 회귀변수 벡터(방문 맥락 + 통제변수) 및 계수 \\
  & $\delta_j$                         & 관광지 고정효과(관광지별 1개) \\ \midrule
\multirow{6}{*}{M2, M3}
  & $i,\,j$                            & 관광지 $j$의 리뷰 $i$ \\
  & $\beta_0$                          & 절편 \\
  & $\mathbf{x}_{ij}^\top\boldsymbol\beta$ & 회귀변수 벡터 및 계수 \\
  & $\delta_j$                         & 관광지 고정효과 \\
  & $\tau_k$                           & 순서형 로짓의 $k$ 범주 임계값 \\
  & $k$                                & 종속변수 범주 (M2: $0\!\dots\!4$;\; M3: $1\!\dots\!4$) \\ \midrule
\multirow{7}{*}{M4}
  & $i,\,j$                            & 관광지 $j$의 리뷰 $i$ \\
  & $\beta_0$                          & 절편 \\
  & $\mathbf{x}_{ij}^\top\boldsymbol\beta$ & 회귀변수 벡터 및 계수 \\
  & $\delta_j$                         & 관광지 고정효과 \\
  & $\mu_{ij}$                         & 반응 수의 기댓값 \\
  & $\gamma$                           & 로그 경과일 탄력성(리뷰 경과일 오프셋) \\
  & $\alpha$                           & 음이항 산포모수($\hat\alpha = 2.44$; 분산 $\approx$ 평균의 37배) \\ \bottomrule
\end{tabular}
\end{table}

\tblVariableMapKO

\tblVIF

\sect{9. 연구결과}
\vspace{-3pt}
\begin{table}[H]\centering
\caption{결과 요약 / 주요 효과 (승산비 및 발생비)}
\footnotesize
\begin{tabular}{lllr}\toprule
\textbf{모형} & \textbf{적합도} & \textbf{변수} & \textbf{OR / IRR} \\ \midrule
\multirow{5}{*}{M1 재방문} & \multirow{5}{*}{\shortstack[l]{AUC 0.773\\(확장 0.785)}}
  & 작성자 총 리뷰 수   & 1.52*** \\
  & & 혼자 방문           & 1.68**  \\
  & & 대기 경험           & 0.51*** \\
  & & 감성 점수           & 0.70*** \\
  & & 여행 그룹           & 0.56*** \\ \midrule
\multirow{4}{*}{M2 키워드 수} & \multirow{4}{*}{\shortstack[l]{RMSE 1.37\\(기준 1.88)}}
  & 텍스트 작성         & 4.02*** \\
  & & 방문 정보 제공      & 2.53*** \\
  & & 재방문자            & 1.60*** \\
  & & 여행 그룹           & 1.73*** \\ \midrule
\multirow{2}{*}{M3 별점} & \multirow{2}{*}{탐색적}
  & 키워드 수           & 1.56*** \\
  & & 재방문자            & 1.37*   \\ \midrule
\multirow{3}{*}{M4 반응 수} & \multirow{3}{*}{AUC 0.815}
  & 팔로워 수           & 1.97*** \\
  & & 텍스트 길이         & 1.64*** \\
  & & 감성 점수           & 0.77**  \\ \bottomrule
\end{tabular}
\end{table}
{\scriptsize * $p<0.05$, ** $p<0.01$, *** $p<0.001$. 비율 $<1$은 부정적 연관을 나타냄. OR/IRR $= e^{\hat\beta}$.}

LLM 기반 분류기로 11{,}164건의 리뷰를 여섯 개 내용 차원(역사·문화, 접근성, 혼잡, 음식,
시각·사진, 서비스·체험)에 라벨링한 측면 수준 분석은 그림을 더 선명하게 한다. 여섯 측면은
관광 문헌 근거로 사전 지정되었으며, LLM 분류 전에 비지도 토픽 추출(LDA 최적 $K=6$, BERTopic
$K=15$)로 보완 검증되었다. 아이와 방문은 세 측면(역사, 접근성, 서비스)의 언급을 1.22--1.55배
높인다. 서비스 마찰은 두 방향으로 분기한다: 대기는 혼잡 언급을 두 배 이상 높이고(OR $2.13$),
사전 예약은 혼잡을 줄이면서(OR $0.69$) 서비스 언급을 두 배 이상 높인다(OR $2.39$). 데이트
목적은 음식 언급을 51\% 상승시킨다. 몰입에 대해서는 작성자 통제 후 어떤 측면 언급도 $p<.05$에
도달하지 않는다: 네이버 플레이스에서 반응은 리뷰 내용이 아닌 작성자 정체성(팔로워, 재방문 여부,
텍스트 길이)에 의해 좌우된다.

\sect{10. 참고문헌}
\begin{enumerate}\small
  \item Bae, E., Chang, M., Park, E., \& Kim, D. (2017). The effect of Hallyu on tourism in Korea. \textit{Journal of Open Innovation: Technology, Market, and Complexity}, 3, 22.
  \item Chang, S. (2021). Food tourism in Korea. \textit{Journal of Vacation Marketing}.
  \item Ko, S., Kang, S., Kang, H., \& Lee, M.~J. (2018). An exploration of foreign tourists' perceptions of Korean food tour. \textit{Asia Pacific Journal of Tourism Research}, 23(8), 833--846.
  \item Karabila, I., Darraz, N., EL-Ansari, A., Alami, N., \& EL~Mallahi, M. (2024). BERT-enhanced sentiment analysis for personalized e-commerce recommendations. \textit{Multimedia Tools and Applications}, 83, 56463--56488.
  \item Lee, M., Jeong, D., \& Kim, A.~J. (2026). The impact of information diversity on online reviews and business performance. \textit{Journal of Consumer Studies}, 37(1). \url{https://doi.org/10.35736/JCS.37.1.1}
  \item Wang, S., Karmakar, S., Wang, F., \& Pei, Y. (2025). Content dissimilarity and online review helpfulness: Contextual insights. \textit{Journal of Business Research}, 187, 115068.
\end{enumerate}

\end{document}
